\documentclass[10pt,a4paper]{amsart}
\usepackage{amsmath,amssymb,amsthm,graphicx}
% Retained minimal preamble for the Stacks Project chapter more-algebra.tex (More on Algebra).
% The chapter file pinned-original-source.tex carries only \input{preamble} on its first line, so no
% preamble text lives inside the chapter itself.  The Stacks Project uses one single shared preamble for
% every chapter; the snapshot stages it identically as sources/stacks-chow/preamble.tex and
% sources/stacks-etale-cohomology/preamble.tex.  The block below is copied mechanically, line for line,
% from sources/stacks-chow/preamble.tex lines 14-18 (the Xy-pic setup) and 154-217 (the native theorem
% declarations together with the mathematical macros).  Nothing was added, removed or redefined; the
% wrapper class and the amsmath/amssymb/amsthm/graphicx base packages are supplied by the assembler.
% Chapter: more-algebra.tex -- More on Algebra.
% For commutative diagrams we use Xy-pic
\usepackage[all]{xy}

% We use 2cell for 2-commutative diagrams.
\xyoption{2cell}

% Theorem environments.
%
\theoremstyle{plain}
\newtheorem{theorem}[subsection]{Theorem}
\newtheorem{proposition}[subsection]{Proposition}
\newtheorem{lemma}[subsection]{Lemma}

\theoremstyle{definition}
\newtheorem{definition}[subsection]{Definition}
\newtheorem{example}[subsection]{Example}
\newtheorem{exercise}[subsection]{Exercise}
\newtheorem{situation}[subsection]{Situation}

\theoremstyle{remark}
\newtheorem{remark}[subsection]{Remark}
\newtheorem{remarks}[subsection]{Remarks}

\numberwithin{equation}{subsection}

% Macros
%
\def\lim{\mathop{\mathrm{lim}}\nolimits}
\def\colim{\mathop{\mathrm{colim}}\nolimits}
\def\Spec{\mathop{\mathrm{Spec}}}
\def\Hom{\mathop{\mathrm{Hom}}\nolimits}
\def\Ext{\mathop{\mathrm{Ext}}\nolimits}
\def\SheafHom{\mathop{\mathcal{H}\!\mathit{om}}\nolimits}
\def\SheafExt{\mathop{\mathcal{E}\!\mathit{xt}}\nolimits}
\def\Sch{\mathit{Sch}}
\def\Mor{\mathop{\mathrm{Mor}}\nolimits}
\def\Ob{\mathop{\mathrm{Ob}}\nolimits}
\def\Sh{\mathop{\mathit{Sh}}\nolimits}
\def\NL{\mathop{N\!L}\nolimits}
\def\CH{\mathop{\mathrm{CH}}\nolimits}
\def\proetale{{pro\text{-}\acute{e}tale}}
\def\etale{{\acute{e}tale}}
\def\QCoh{\mathit{QCoh}}
\def\Ker{\mathop{\mathrm{Ker}}}
\def\Im{\mathop{\mathrm{Im}}}
\def\Coker{\mathop{\mathrm{Coker}}}
\def\Coim{\mathop{\mathrm{Coim}}}

% Boxtimes
%
\DeclareMathSymbol{\boxtimes}{\mathbin}{AMSa}{"02}

%
% Macros for moduli stacks/spaces
%
\def\QCohstack{\mathcal{QC}\!\mathit{oh}}
\def\Cohstack{\mathcal{C}\!\mathit{oh}}
\def\Spacesstack{\mathcal{S}\!\mathit{paces}}
\def\Quotfunctor{\mathrm{Quot}}
\def\Hilbfunctor{\mathrm{Hilb}}
\def\Curvesstack{\mathcal{C}\!\mathit{urves}}
\def\Polarizedstack{\mathcal{P}\!\mathit{olarized}}
\def\Complexesstack{\mathcal{C}\!\mathit{omplexes}}
% \Pic is the operator that assigns to X its picard group, usage \Pic(X)
% \Picardstack_{X/B} denotes the Picard stack of X over B
% \Picardfunctor_{X/B} denotes the Picard functor of X over B
\def\Pic{\mathop{\mathrm{Pic}}\nolimits}
\def\Picardstack{\mathcal{P}\!\mathit{ic}}
\def\Picardfunctor{\mathrm{Pic}}
\def\Deformationcategory{\mathcal{D}\!\mathit{ef}}

% --- appended 2026-10-12 (Chief): the project own definitions of the reference/slogan/history
% display environments, copied mechanically from the shared Stacks preamble lines 8-12
% (\usepackage{verbatim} plus the three \newenvironment lines, which the chapter preamble keeps
% verbatim from the project preamble). Nothing added, removed or redefined: these environments
% carry no mathematical content and the project itself typesets them as comments. Without these
% three lines a counted statement whose own text uses them fails with Environment slogan undefined.
\usepackage{verbatim}
\newenvironment{reference}{\comment}{\endcomment}
\newenvironment{slogan}{\comment}{\endcomment}
\newenvironment{history}{\comment}{\endcomment}
\usepackage{cleveref}
\crefname{originallemma}{original lemma}{original lemmas}
\Crefname{originallemma}{Original lemma}{Original lemmas}
\makeatletter
\newlabel{algebra-lemma-colimit-category-fp-modules}{{05N7}{source}{}{sourcepointer-1}{}}
\newlabel{algebra-lemma-colimit-category-fp-modules@cref}{{[originallemma][0][]05N7}{[1][0][]source}{}{}{}}
\makeatother
\crefname{originallemma}{original lemma}{original lemmas}
\Crefname{originallemma}{Original lemma}{Original lemmas}
\makeatletter
\newlabel{algebra-lemma-localization-colimit}{{00CR}{source}{}{sourcepointer-2}{}}
\newlabel{algebra-lemma-localization-colimit@cref}{{[originallemma][0][]00CR}{[1][0][]source}{}{}{}}
\makeatother
\crefname{originallemma}{original lemma}{original lemmas}
\Crefname{originallemma}{Original lemma}{Original lemmas}
\makeatletter
\newlabel{algebra-lemma-universally-exact-split}{{058L}{source}{}{sourcepointer-3}{}}
\newlabel{algebra-lemma-universally-exact-split@cref}{{[originallemma][0][]058L}{[1][0][]source}{}{}{}}
\makeatother
\crefname{originallemma}{original lemma}{original lemmas}
\Crefname{originallemma}{Original lemma}{Original lemmas}
\makeatletter
\newlabel{algebra-lemma-universally-injective-check-stalks}{{05CL}{source}{}{sourcepointer-4}{}}
\newlabel{algebra-lemma-universally-injective-check-stalks@cref}{{[originallemma][0][]05CL}{[1][0][]source}{}{}{}}
\makeatother
\crefname{originallemma}{original lemma}{original lemmas}
\Crefname{originallemma}{Original lemma}{Original lemmas}
\makeatletter
\newlabel{lemma-which-elements-split}{{38731}{source}{}{sourcepointer-5}{}}
\newlabel{lemma-which-elements-split@cref}{{[originallemma][0][]38731}{[1][0][]source}{}{}{}}
\makeatother
\crefname{originalproposition}{original proposition}{original propositions}
\Crefname{originalproposition}{Original proposition}{Original propositions}
\makeatletter
\newlabel{proposition-splitting}{{38809}{source}{}{sourcepointer-6}{}}
\newlabel{proposition-splitting@cref}{{[originalproposition][0][]38809}{[1][0][]source}{}{}{}}
\makeatother
\title{Over a ring with finite-dimensional Noetherian maximum spectrum a finitely presented module with a large free summand locally splits off a free summand globally}
\author{The Stacks Project Authors}
\date{stacks-project:more-algebra.tex:local-snapshot-20261012; GNU FDL 1.2 or later}
\begin{document}
\maketitle
\setcounter{subsection}{1}
\noindent\textit{Editorial note.} All delivered mathematical text is the author's own, copied from the pinned Stacks Project chapter file more-algebra.tex; no statement or proof step was rewritten, reordered or removed.\par\medskip
\noindent\textit{Editorial note.} The chapter file pinned-original-source.tex carries only \textbackslash{}input\{preamble\} on its first line, so no preamble text lives inside the chapter. The retained file sources/stacks-more-algebra/retained-preamble.tex is a mechanical, line-for-line copy of the project's single shared preamble as staged in the snapshot, together with the project's own definitions of the reference/slogan/history display environments, which the project itself typesets as comments; nothing was added, removed or redefined by this card and no package was added.\par\medskip
\noindent\textit{Editorial note.} The source URL is the source-level Stacks Project URL, as required when more than one card is produced from one source; the per-claim tag is carried by claim\_key.\par\medskip
\noindent\textit{Editorial note.} The counted statement span includes the author's own \textbackslash{}begin\{reference\} block citing the published theorem; it is delivered verbatim and produces no unresolved citation. The author writes the enumeration of the specialization of the splitting proposition in full ('with F = empty, h = 0 and no s\_i, n = 0 and no x\_i, v\_i, and k = d + 1'); that sentence is reproduced unchanged. The six references inside the counted proof span are recorded as non-load-bearing navigation pointers.\par\medskip
\noindent\textit{Allowed standard prerequisites (editorial).}\begin{itemize}
\item Spreading out a local direct summand to a principal open, via the colimit descriptions of localization and of finitely presented modules. used at proof lines 38933-38938. Converts the local rank hypothesis into the global dimension statement about V(x).
\item The splitting proposition producing an element whose nonvanishing locus is controlled by a codimension bound. invoked at proof lines 38941-38945. Is the main engine of the proof, applied with k = d + 1.
\item Universal injectivity checked at stalks, and universally injective maps are split injective. invoked at proof lines 38947-38951. Upgrades the pointwise summand statement into a genuine direct sum decomposition of M.
\end{itemize}\medskip
\ifdefined\hypertarget\hypertarget{sourcepointer-1}{}\fi
\noindent\textit{Source reference.} Original lemma 05N7: focused-seeds/AGP/sources/stacks-snapshot-20261011/tags.txt, source lines 7296--7296. Cross-chapter navigation pointer into the Stacks Project; the standard fact it supplies is applied inside the counted proof span and the argument resting on it stays there.\par\medskip
\ifdefined\hypertarget\hypertarget{sourcepointer-2}{}\fi
\noindent\textit{Source reference.} Original lemma 00CR: focused-seeds/AGP/sources/stacks-snapshot-20261011/tags.txt, source lines 471--471. Cross-chapter navigation pointer into the Stacks Project; the standard fact it supplies is applied inside the counted proof span and the argument resting on it stays there.\par\medskip
\ifdefined\hypertarget\hypertarget{sourcepointer-3}{}\fi
\noindent\textit{Source reference.} Original lemma 058L: focused-seeds/AGP/sources/stacks-snapshot-20261011/tags.txt, source lines 6784--6784. Cross-chapter navigation pointer into the Stacks Project; the standard fact it supplies is applied inside the counted proof span and the argument resting on it stays there.\par\medskip
\ifdefined\hypertarget\hypertarget{sourcepointer-4}{}\fi
\noindent\textit{Source reference.} Original lemma 05CL: focused-seeds/AGP/sources/stacks-snapshot-20261011/tags.txt, source lines 6924--6924. Cross-chapter navigation pointer into the Stacks Project; the standard fact it supplies is applied inside the counted proof span and the argument resting on it stays there.\par\medskip
\ifdefined\hypertarget\hypertarget{sourcepointer-5}{}\fi
\noindent\textit{Source reference.} Original lemma 38731: sources/stacks-more-algebra/pinned-original-source.tex, source lines 38731--38731. Same-chapter navigation pointer; the statement it names is standard for the section and the argument that consumes it stays inside the counted proof span.\par\medskip
\ifdefined\hypertarget\hypertarget{sourcepointer-6}{}\fi
\noindent\textit{Source reference.} Original proposition 38809: sources/stacks-more-algebra/pinned-original-source.tex, source lines 38809--38809. Same-chapter navigation pointer; the statement it names is standard for the section and the argument that consumes it stays inside the counted proof span.\par\medskip
\begin{theorem}
\label{theorem-splitting}
\begin{reference}
\cite[Theorem 1]{Serre-projective}
\end{reference}
Let $R$ be a ring whose max spectrum $\Omega \subset \Spec(R)$
is a Noetherian topological space of dimension $d < \infty$.
Let $M$ be a finitely presented $R$-module such that
for all $\mathfrak m \in \Omega$ the $R_\mathfrak m$-module
$M_\mathfrak m$ has a free direct summand of rank $> d$.
Then $M \cong R \oplus M'$.
\end{theorem}

\begin{proof}
For $\mathfrak m \in \Omega$ suppose that $R_\mathfrak m^{\oplus r}$
is a direct summand of $M_\mathfrak m$. Then by Algebra, Lemmas
\ref{algebra-lemma-localization-colimit} and
\ref{algebra-lemma-colimit-category-fp-modules}
we see that $R_f^{\oplus r}$ is a direct summand of $M_f$
for some $f \in R$, $f \not \in \mathfrak m$. Hence
the assumption means that $\dim V(x) > d$ for all $x \in \Omega$
where $V(x)$ is as in Lemma \ref{lemma-which-elements-split}.
By Proposition \ref{proposition-splitting} applied with $F = \emptyset$,
$h = 0$ and no $s_i$, $n = 0$ and no $x_i, v_i$, and $k = d + 1$
we find an $s \in M$ and $F' \subset \Omega$
such that every irreducible component of $F'$ has codimension $\geq d + 1$
and $Z(s) \subset F'$. Since $d = \dim(\Omega)$ this forces
$F' = \emptyset$. Hence $s : R \to M$ is the inclusion of a
direct summand at all maximal ideals. It follows that $s$ is universally
injective, see Algebra, Lemma
\ref{algebra-lemma-universally-injective-check-stalks}. Then
$s$ is split injective by
Algebra, Lemma \ref{algebra-lemma-universally-exact-split}.
\end{proof}

\end{document}
